\section{The bitopological settling site}

\begin{proposition}[The FOUR has a finite Priestley dual]\label{thm:four-priestley-dual}\lean{Obsmetrologia.Stone.four_priestley_dual}\leanok\uses{thm:four-object, thm:four-bilattice}
The four charges are the clopen up-sets of the two-point pairwise-Stone dual.
This is the local topological carrier over which the settling site is built.
\end{proposition}

\begin{theorem}[The pairwise-Stone topology recovers the order]\label{thm:pairwise-stone-order-recovery}\lean{Obsmetrologia.Stone.order_recovered_from_tauPlus}\leanok\uses{thm:four-priestley-dual}
The bitopological pair recovers the Priestley order: the order is exactly the
relation preserving all positive clopen observations.
\end{theorem}

\begin{theorem}[The bitopological settling site is Grothendieck]\label{thm:bitopological-site-grothendieck}\lean{Obsmetrologia.Site.bitopological_site_grothendieck}\leanok\uses{thm:pairwise-stone-order-recovery, thm:closure-regime-carrier}
The objects are predicate, stratum, and clopen data; the covering relation is a
join-cover in the pairwise-Stone clopen fiber.  The clopen algebra, inclusion
order, temporal preorder, and Grothendieck basis axioms are all formalized.
\end{theorem}

\begin{theorem}[Covers pull back to clopen overlaps]\label{thm:covers-pullback}\lean{Obsmetrologia.Site.covers_pullback}\leanok\uses{thm:bitopological-site-grothendieck}
Every cover of a clopen pulls back along a clopen inclusion: if a family covers
\(U\) and \(V \subseteq U\), then the family of intersections \(V \cap U_i\)
covers \(V\).  This is the site-level base-change law used to identify pairwise
overlaps as genuine local sites rather than extra data.
\end{theorem}

\begin{proposition}[Claim site points]\label{thm:claim-site-point}\lean{Obsmetrologia.Claim.SitePoint}\leanok\uses{thm:claim-point, thm:bitopological-site-grothendieck}
A claim site point is a site object over a claim point: claim, stratum, and
clopen support. The point is generated by predication; the site data is supplied
by the settling site. Covers and descent remain the structure of the site, not
extra fields on the claim itself.
\end{proposition}

\begin{theorem}[Claim fiber morphisms retain support]\label{thm:claim-fiberhom-support}\lean{Obsmetrologia.Claim.fiberHom_support}\leanok\uses{thm:claim-site-point}
Fiber morphisms between claim site points are typed morphisms, not mere
propositions: they retain the source clopen support as morphism data while
carrying the fixed-claim, fixed-stratum, and inclusion witnesses.
\end{theorem}

\begin{proposition}[Claim term contexts]\label{thm:claim-term-context}\lean{Obsmetrologia.Claim.TermContext}\leanok\uses{thm:claim-point}
A term context is a generated term carrier together with incidence against
claim points.  The term carrier is data of the context, not a new primitive:
different generated syntaxes may supply the carrier once their incidence against
claim points is proved.
\end{proposition}

\begin{proposition}[Claim concepts]\label{thm:claim-concept}\lean{Obsmetrologia.Claim.Concept}\leanok\uses{thm:claim-term-context, thm:fca-closure-laws}
A claim concept is the existing FCA concept specialized to claim points and the
term carrier of a term context.  Its extent and intent retain the FCA derivation
witnesses; in particular, the extent is closed under the context's derivation
closure.
\end{proposition}

\begin{proposition}[Claim topics]\label{thm:claim-topic}\lean{Obsmetrologia.Claim.Topic}\leanok\uses{thm:claim-concept, thm:claim-site-point}
A topic is a concept localized at a claim site point together with an explicit
cover of that point's clopen support.  It therefore carries the concept data,
the membership of the point's claim in the concept extent, and the site cover
data needed for descent.
\end{proposition}

\begin{theorem}[Topic covers are fiber-local]\label{thm:claim-topic-cover-fiber-local}\lean{Obsmetrologia.Claim.Topic.cover_is_fiber_local}\leanok\uses{thm:claim-topic, thm:no-temporal-covers}
The cover carried by a topic is still the site's clopen cover: every member lies
below the topic point's support, and the family joins back to that support.
\end{theorem}

\begin{theorem}[Topics restrict along claim fibers]\label{thm:claim-topic-restrict-cover}\lean{Obsmetrologia.Claim.Topic.restrict_cover_is_fiber_local}\leanok\uses{thm:claim-topic-cover-fiber-local, thm:covers-pullback, thm:claim-fiberhom-support}
A topic restricts along a fixed-claim, fixed-stratum clopen inclusion by pulling
its cover back to the source support.  The concept and extent witness remain the
same carried concept data; only the site point and cover are base-changed.
\end{theorem}

\begin{theorem}[Topic covers generate local topic pieces]\label{thm:claim-topic-cover-restriction}\lean{Obsmetrologia.Claim.Topic.coverRestriction_cover_is_fiber_local}\leanok\uses{thm:claim-topic-restrict-cover}
Each member of a topic's own cover determines a claim site point over that
member and a restricted local topic there.  Thus the cover is not inert
metadata: it generates the local pieces of the topic.
\end{theorem}

\begin{theorem}[Topic local pieces agree on overlaps]\label{thm:claim-topic-cover-overlap-concepts}\lean{Obsmetrologia.Claim.Topic.coverRestriction_overlap_concepts_agree}\leanok\uses{thm:claim-topic-cover-restriction}
The local topics generated by two members of one topic cover carry the same
concept when both are further restricted to their clopen overlap.  This is the
first compatibility law for topic descent.
\end{theorem}

\begin{proposition}[Term profiles]\label{thm:claim-term-profile}\lean{Obsmetrologia.Claim.TermContext.profile}\leanok\uses{thm:claim-term-context}
A term in a term context has a direct claim profile: the claims incident with
that term.  This is the first readable surface of a term; the bare carrier
element alone carries no such profile.
\end{proposition}

\begin{theorem}[Terms generate closed concepts]\label{thm:claim-term-generated-concept}\lean{Obsmetrologia.Claim.TermContext.generatedConcept_extent_closed}\leanok\uses{thm:claim-term-profile, thm:claim-concept}
The direct claim profile of a term closes to an FCA concept, and that generated
concept has a closed extent under the context's derivation closure.
\end{theorem}

\begin{theorem}[Generated terms lie in generated intents]\label{thm:claim-term-generated-intent}\lean{Obsmetrologia.Claim.TermContext.generatedConcept_in_intent}\leanok\uses{thm:claim-term-generated-concept}
The term whose direct profile generates a concept lies in that concept's intent:
every claim in the direct profile is incident with the generating term by
definition.
\end{theorem}

\begin{proposition}[Term-context transport]\label{thm:claim-term-context-hom}\lean{Obsmetrologia.Claim.TermContext.Hom}\leanok\uses{thm:claim-term-context}
A transport between term contexts is a claim map and a term map with an
incidence-preservation witness.  This is the data needed before one may speak
of the same reading moving across changed contexts.
\end{proposition}

\begin{theorem}[Term profiles transport]\label{thm:claim-term-profile-transport}\lean{Obsmetrologia.Claim.TermContext.Hom.profile_preserved}\leanok\uses{thm:claim-term-context-hom, thm:claim-term-profile}
A term-context transport carries a term's direct profile into the transported
term's direct profile.  The proof uses the incidence-preservation witness rather
than identifying the two contexts definitionally.
\end{theorem}

\begin{theorem}[Valuation support pulls back along term transport]\label{thm:claim-term-profile-valuation-pullback}\lean{Obsmetrologia.Claim.TermContext.Hom.pullback_supports_profile}\leanok\uses{thm:claim-term-profile-transport, thm:reindex}
If a target valuation truth-supports the transported profile of a term, then
the valuation pulled back along the term-context transport supports the source
profile.  This uses the existing valuation reindexing law rather than a
separate claim-by-claim convention.
\end{theorem}

\begin{theorem}[Claim-fixing transports have identity valuation pullback]\label{thm:claim-term-valuation-pullback-claim-fixing}\lean{Obsmetrologia.Claim.TermContext.Hom.pullbackValuation_eq_self_of_claimMap_eq}\leanok\uses{thm:claim-term-context-hom, thm:reindex}
If a term-context transport fixes every claim, then pulling a valuation back
along that transport recovers the original valuation.
\end{theorem}

\begin{theorem}[Valuation pullback commutes with closure]\label{thm:claim-term-valuation-pullback-closure}\lean{Obsmetrologia.Claim.TermContext.Hom.pullbackValuation_liftOp}\leanok\uses{thm:claim-term-profile-valuation-pullback, thm:reindex, thm:closure-operator}
Pulling a valuation back along a term-context transport commutes with pointwise
valuation closure.  The proof is the existing \(C\)-naturality of lifted closure
read at the transport's claim map.
\end{theorem}

\begin{proposition}[Term readings]\label{thm:claim-term-reading}\lean{Obsmetrologia.Claim.TermReading}\leanok\uses{thm:claim-topic}
A term reading is a term carried by the closed intent of a local topic.  It
therefore combines the term carrier, the FCA intent, the localized claim-site
point, and the topic cover in one proof-carrying datum.
\end{proposition}

\begin{theorem}[Term readings restrict along claim fibers]\label{thm:claim-term-reading-restrict-cover}\lean{Obsmetrologia.Claim.TermReading.restrict_cover_is_fiber_local}\leanok\uses{thm:claim-term-reading, thm:claim-topic-restrict-cover}
A term reading restricts by restricting its topic along the claim-site fiber
morphism while retaining the same term and intent witness.  The restricted
reading therefore carries the source point's fiber-local cover.
\end{theorem}

\begin{theorem}[Term reading covers generate local readings]\label{thm:claim-term-reading-cover-restriction}\lean{Obsmetrologia.Claim.TermReading.coverRestriction_cover_is_fiber_local}\leanok\uses{thm:claim-term-reading-restrict-cover, thm:claim-topic-cover-restriction}
Each member of a term reading's topic cover determines a local term reading
whose topic is the corresponding restricted topic and whose term remains the
same carried term.
\end{theorem}

\begin{theorem}[Term reading local pieces agree on overlaps]\label{thm:claim-term-reading-cover-overlap-terms}\lean{Obsmetrologia.Claim.TermReading.coverRestriction_overlap_terms_agree}\leanok\uses{thm:claim-term-reading-cover-restriction, thm:claim-topic-cover-overlap-concepts}
The local term readings generated by two cover members carry the same term when
restricted again to their overlap.  The term is not rediscovered locally; it is
the parent reading's carried term.
\end{theorem}

\begin{proposition}[Valuated term readings]\label{thm:claim-valuated-term-reading}\lean{Obsmetrologia.Claim.ValuatedTermReading}\leanok\uses{thm:claim-term-reading, thm:claim-term-profile, thm:valuation}
A valuated term reading is a term reading together with a valuation and a
witness that every claim in the term's direct profile fires truthwise under that
valuation.  It is a carried reading, not a bare proposition about a term.
\end{proposition}

\begin{theorem}[Valuated readings restrict with support]\label{thm:claim-valuated-reading-restrict-support}\lean{Obsmetrologia.Claim.ValuatedTermReading.restrict_supports_profile}\leanok\uses{thm:claim-valuated-term-reading, thm:claim-term-reading-restrict-cover}
A valuated term reading restricts along the topic restriction while retaining
the valuation and direct-profile support witness.  Support is therefore not
recomputed from a proposition; it is carried through the restricted reading.
\end{theorem}

\begin{theorem}[Valuated reading restrictions remain fiber-local]\label{thm:claim-valuated-reading-restrict-cover}\lean{Obsmetrologia.Claim.ValuatedTermReading.restrict_cover_is_fiber_local}\leanok\uses{thm:claim-valuated-reading-restrict-support, thm:claim-term-reading-restrict-cover}
The restricted valuated reading still carries a cover local to the source claim
site point, because its underlying term reading restricts by topic base-change.
\end{theorem}

\begin{theorem}[Valuated reading covers generate supported local readings]\label{thm:claim-valuated-reading-cover-restriction-support}\lean{Obsmetrologia.Claim.ValuatedTermReading.coverRestriction_supports_profile}\leanok\uses{thm:claim-valuated-reading-restrict-support, thm:claim-term-reading-cover-restriction}
Each member of a valuated reading's topic cover determines a local valuated
reading that retains the original valuation and direct-profile support witness.
\end{theorem}

\begin{theorem}[Cover-restricted valuated readings remain fiber-local]\label{thm:claim-valuated-reading-cover-restriction-cover}\lean{Obsmetrologia.Claim.ValuatedTermReading.coverRestriction_cover_is_fiber_local}\leanok\uses{thm:claim-valuated-reading-cover-restriction-support, thm:claim-valuated-reading-restrict-cover}
The local valuated reading generated by a cover member carries a fiber-local
cover over that member.  This is the cover-indexed version of valuated reading
restriction.
\end{theorem}

\begin{theorem}[Valuated reading local pieces agree on overlaps]\label{thm:claim-valuated-reading-cover-overlap-valuations}\lean{Obsmetrologia.Claim.ValuatedTermReading.coverRestriction_overlap_valuations_agree}\leanok\uses{thm:claim-valuated-reading-cover-restriction-cover, thm:claim-term-reading-cover-overlap-terms}
The local valuated readings generated by two cover members carry the same
valuation when restricted again to their overlap.  Valuation support is therefore
transported as part of one parent reading rather than chosen independently on
patches.
\end{theorem}

\begin{proposition}[Clopens read valuations as cell patches]\label{thm:claim-valuated-reading-clopen-patch}\lean{Obsmetrologia.Claim.ValuatedTermReading.clopenPatch}\leanok\uses{thm:four-priestley-dual, thm:claim-valuated-term-reading}
A pairwise-Stone clopen has positive and negative evidence slots, and a valuation
assigns each predication cell a charge with the same two slots.  Reading a clopen
through a valuation therefore gives a cell patch: the cell is admitted when the
clopen admits at least one slot present in the cell's charge.
\end{proposition}

\begin{theorem}[Valuation-read clopen joins are patch unions]\label{thm:claim-valuated-reading-clopen-patch-join}\lean{Obsmetrologia.Claim.ValuatedTermReading.clopenPatch_join}\leanok\uses{thm:claim-valuated-reading-clopen-patch, thm:bitopological-site-grothendieck}
Reading a clopen join through a valuation is the Boolean union of the two
valuation-read patches.  This is the byte-level bridge from the site clopen
algebra to sheaf-style cell masks.
\end{theorem}

\begin{theorem}[Finite clopen covers generate patch members]\label{thm:claim-valuated-reading-clopen-patch-joinlist}\lean{Obsmetrologia.Claim.ValuatedTermReading.clopenPatch_joinList_true}\leanok\uses{thm:claim-valuated-reading-clopen-patch-join}
If a cell lies in the valuation-read patch of a finite clopen join, then one
member of that clopen family already admits the cell.  The proof uses the finite
pairwise-Stone slot algebra, not a new cover convention.
\end{theorem}

\begin{proposition}[Valuated topic covers form relative patch families]\label{thm:claim-valuated-reading-topic-patch-family}\lean{Obsmetrologia.Claim.ValuatedTermReading.topicPatchFamily}\leanok\uses{thm:claim-valuated-reading-clopen-patch-joinlist, thm:claim-valuated-reading-restrict-cover}
A valuated reading's topic cover induces a relative patch family over the
valuation-read support of the topic point.  It covers that support mask; it does
not claim to cover every predication cell.
\end{proposition}

\begin{proposition}[Full-support topic patches give sheaf covers]\label{thm:claim-valuated-reading-topic-sheaf-cover}\lean{Obsmetrologia.Claim.ValuatedTermReading.topicSheafCover}\leanok\uses{thm:claim-valuated-reading-topic-patch-family}
The relative patch family becomes a global settling sheaf cover only with an
explicit hypothesis that the valuation-read topic support contains every
predication cell.  This records exactly which hypothesis the sheaf boundary
depends on.
\end{proposition}

\begin{proposition}[Intensional frames]\label{thm:claim-intensional-frame}\lean{Obsmetrologia.Claim.IntensionalFrame}\leanok\uses{thm:claim-term-context}
An intensional frame is a family of term contexts indexed by scenario-like
positions.  This is the place where context variation becomes data rather than
a prose assumption.
\end{proposition}

\begin{theorem}[Primary readings close pointwise]\label{thm:claim-primary-concept-closed}\lean{Obsmetrologia.Claim.IntensionalFrame.primaryConcept_extent_closed}\leanok\uses{thm:claim-intensional-frame, thm:claim-term-generated-concept}
A primary concept reading chooses a local term at each index and closes that
term inside that index's term context.  Each such pointwise concept has a closed
extent.
\end{theorem}

\begin{theorem}[Primary readings localize to topics]\label{thm:claim-primary-topic-cover}\lean{Obsmetrologia.Claim.IntensionalFrame.primaryTopic_cover_is_fiber_local}\leanok\uses{thm:claim-primary-concept-closed, thm:claim-topic}
Once a site point is witnessed in a primary concept's extent, that primary
concept forms a topic at the point, and the resulting topic cover is fiber-local.
\end{theorem}

\begin{theorem}[Primary topics carry term readings]\label{thm:claim-primary-reading-cover}\lean{Obsmetrologia.Claim.IntensionalFrame.primaryReading_cover_is_fiber_local}\leanok\uses{thm:claim-primary-topic-cover, thm:claim-term-generated-intent, thm:claim-term-reading}
The primary topic also carries the locally chosen term as a term reading, because
the generating term lies in the generated concept's intent.  Its cover remains
the topic's fiber-local cover.
\end{theorem}

\begin{theorem}[Primary readings carry valuation support]\label{thm:claim-primary-valuated-reading-cover}\lean{Obsmetrologia.Claim.IntensionalFrame.primaryValuatedReading_cover_is_fiber_local}\leanok\uses{thm:claim-primary-reading-cover, thm:claim-valuated-term-reading}
If a valuation truth-supports the primary term's direct profile, the primary
reading becomes a valuated term reading.  Its cover remains the underlying
fiber-local topic cover.
\end{theorem}

\begin{proposition}[Primary readings induce relative patch families]\label{thm:claim-primary-patch-family}\lean{Obsmetrologia.Claim.IntensionalFrame.primaryPatchFamily}\leanok\uses{thm:claim-primary-valuated-reading-cover, thm:claim-valuated-reading-topic-patch-family}
A primary valuated reading's topic cover reads through its valuation as a
relative patch family over the primary topic support.  This is primary intension
composed through term choice, concept closure, site localization, valuation
support, and the patch-family bridge.
\end{proposition}

\begin{theorem}[Primary patch families cover their support]\label{thm:claim-primary-patch-family-covers-support}\lean{Obsmetrologia.Claim.IntensionalFrame.primaryPatchFamily_covers_support}\leanok\uses{thm:claim-primary-patch-family}
The primary patch family covers exactly its displayed support mask.  No global
sheaf cover is asserted until full support is supplied.
\end{theorem}

\begin{proposition}[Primary patch families promote under full support]\label{thm:claim-primary-sheaf-cover}\lean{Obsmetrologia.Claim.IntensionalFrame.primarySheafCover}\leanok\uses{thm:claim-primary-patch-family-covers-support, thm:claim-valuated-reading-topic-sheaf-cover}
A primary patch family becomes a settling sheaf cover when the primary topic's
valuation-read support contains every predication cell.
\end{proposition}

\begin{theorem}[Secondary anchors fix the actual concept]\label{thm:claim-secondary-actual-concept}\lean{Obsmetrologia.Claim.IntensionalFrame.secondaryConcept_actual}\leanok\uses{thm:claim-intensional-frame, thm:claim-term-generated-concept}
At the actual index, an anchored secondary reading closes the actual term itself:
the actual transport fixes the anchored term, so the secondary concept there is
the generated concept of the actual term.
\end{theorem}

\begin{theorem}[Secondary readings preserve actual profiles]\label{thm:claim-secondary-profile-preserved}\lean{Obsmetrologia.Claim.IntensionalFrame.secondary_profile_preserved}\leanok\uses{thm:claim-intensional-frame, thm:claim-term-profile-transport}
A secondary concept reading fixes one actual term and transports it into each
context before closing it.  The actual term's direct profile is preserved along
the chosen term-context transport.
\end{theorem}

\begin{theorem}[Secondary actual profiles remain actual]\label{thm:claim-secondary-profile-preserved-actual}\lean{Obsmetrologia.Claim.IntensionalFrame.secondary_profile_preserved_at_actual}\leanok\uses{thm:claim-secondary-profile-preserved, thm:claim-secondary-actual-concept}
At the actual index, the transported profile is still read at the same claim:
the actual transport fixes claims as well as the anchored term.
\end{theorem}

\begin{proposition}[Secondary anchors induce primary term families]\label{thm:claim-anchor-induced-term-family}\lean{Obsmetrologia.Claim.IntensionalFrame.AnchoredTerm.termFamily}\leanok\uses{thm:claim-intensional-frame}
An anchored secondary term determines a primary term family by transporting the
actual term into every context.  This is the exact construction by which a
secondary intension can be read as a primary intension, not a definitional
renaming of the two notions.
\end{proposition}

\begin{theorem}[Anchor-induced families recover the actual term]\label{thm:claim-anchor-induced-term-family-actual}\lean{Obsmetrologia.Claim.IntensionalFrame.AnchoredTerm.termFamily_actual}\leanok\uses{thm:claim-anchor-induced-term-family}
At the actual index, the induced primary family returns the anchored actual
term, using the anchor's actual-transport witness.
\end{theorem}

\begin{theorem}[Actual anchor pullback recovers valuations]\label{thm:claim-anchor-actual-pullback-valuation}\lean{Obsmetrologia.Claim.IntensionalFrame.AnchoredTerm.actual_pullbackValuation}\leanok\uses{thm:claim-anchor-induced-term-family-actual, thm:claim-term-valuation-pullback-claim-fixing}
At an anchor's actual index, valuation pullback along the actual transport is
the original valuation, because the anchor carries a claim-fixing witness.
\end{theorem}

\begin{theorem}[Actual support transports from the anchored term]\label{thm:claim-anchor-actual-supports-transport}\lean{Obsmetrologia.Claim.IntensionalFrame.AnchoredTerm.actual_supports_transport}\leanok\uses{thm:claim-anchor-induced-term-family-actual, thm:claim-secondary-profile-preserved-actual}
If a valuation supports the anchored actual term, it also supports the term
obtained by applying the actual transport to that term.
\end{theorem}

\begin{theorem}[Secondary concepts are induced primary concepts]\label{thm:claim-primary-concept-from-anchor}\lean{Obsmetrologia.Claim.IntensionalFrame.primaryConcept_from_anchor}\leanok\uses{thm:claim-anchor-induced-term-family, thm:claim-secondary-actual-concept}
The secondary concept at any index is exactly the primary concept of the
anchor-induced term family at that index.
\end{theorem}

\begin{theorem}[Secondary readings are induced primary readings]\label{thm:claim-primary-reading-from-anchor}\lean{Obsmetrologia.Claim.IntensionalFrame.primaryReading_from_anchor}\leanok\uses{thm:claim-primary-concept-from-anchor, thm:claim-secondary-reading-cover}
After localization at a witnessed site point, the secondary term reading is the
primary term reading of the anchor-induced family.
\end{theorem}

\begin{theorem}[Secondary patch families are induced primary patch families]\label{thm:claim-primary-patch-family-from-anchor}\lean{Obsmetrologia.Claim.IntensionalFrame.primaryPatchFamily_from_anchor}\leanok\uses{thm:claim-primary-reading-from-anchor, thm:claim-secondary-patch-family}
After valuation support and the topic-to-patch bridge, the secondary patch
family is still the primary patch family of the anchor-induced term family.
\end{theorem}

\begin{theorem}[Secondary sheaf covers are induced primary sheaf covers]\label{thm:claim-primary-sheaf-cover-from-anchor}\lean{Obsmetrologia.Claim.IntensionalFrame.primarySheafCover_from_anchor}\leanok\uses{thm:claim-primary-patch-family-from-anchor, thm:claim-secondary-sheaf-cover}
Under the same full-support hypothesis, a secondary sheaf cover is the primary
sheaf cover obtained from the anchor-induced term family.
\end{theorem}

\begin{theorem}[Secondary readings localize to topics]\label{thm:claim-secondary-topic-cover}\lean{Obsmetrologia.Claim.IntensionalFrame.secondaryTopic_cover_is_fiber_local}\leanok\uses{thm:claim-secondary-profile-preserved, thm:claim-topic}
Once a site point is witnessed in a secondary concept's extent, that transported
actual concept forms a topic at the point, and the resulting topic cover is
fiber-local.
\end{theorem}

\begin{theorem}[Secondary topics carry term readings]\label{thm:claim-secondary-reading-cover}\lean{Obsmetrologia.Claim.IntensionalFrame.secondaryReading_cover_is_fiber_local}\leanok\uses{thm:claim-secondary-topic-cover, thm:claim-term-generated-intent, thm:claim-term-reading}
The secondary topic also carries the transported actual term as a term reading,
because that transported term lies in the generated secondary concept's intent.
Its cover remains the topic's fiber-local cover.
\end{theorem}

\begin{theorem}[Secondary readings carry valuation support]\label{thm:claim-secondary-valuated-reading-cover}\lean{Obsmetrologia.Claim.IntensionalFrame.secondaryValuatedReading_cover_is_fiber_local}\leanok\uses{thm:claim-secondary-reading-cover, thm:claim-valuated-term-reading}
If a valuation truth-supports the transported actual term's direct profile in a
target context, the secondary reading in that context becomes a valuated term
reading.  Its cover remains the underlying fiber-local topic cover.
\end{theorem}

\begin{proposition}[Secondary readings induce relative patch families]\label{thm:claim-secondary-patch-family}\lean{Obsmetrologia.Claim.IntensionalFrame.secondaryPatchFamily}\leanok\uses{thm:claim-secondary-valuated-reading-cover, thm:claim-valuated-reading-topic-patch-family}
A secondary valuated reading's transported actual term, closed concept, localized
topic, valuation support, and topic cover compose into a relative patch family.
\end{proposition}

\begin{theorem}[Secondary patch families cover their support]\label{thm:claim-secondary-patch-family-covers-support}\lean{Obsmetrologia.Claim.IntensionalFrame.secondaryPatchFamily_covers_support}\leanok\uses{thm:claim-secondary-patch-family}
The secondary patch family covers its displayed valuation-read support mask,
not an unstated global cell base.
\end{theorem}

\begin{proposition}[Secondary patch families promote under full support]\label{thm:claim-secondary-sheaf-cover}\lean{Obsmetrologia.Claim.IntensionalFrame.secondarySheafCover}\leanok\uses{thm:claim-secondary-patch-family-covers-support, thm:claim-valuated-reading-topic-sheaf-cover}
A secondary patch family becomes a settling sheaf cover only when the secondary
topic's valuation-read support is global.
\end{proposition}

\begin{theorem}[Secondary valuation support pulls back to the actual reading]\label{thm:claim-secondary-pullback-valuated-reading-cover}\lean{Obsmetrologia.Claim.IntensionalFrame.secondaryActualPullbackReading_cover_is_fiber_local}\leanok\uses{thm:claim-secondary-valuated-reading-cover, thm:claim-term-profile-valuation-pullback, thm:claim-secondary-actual-concept}
Target valuation support for the transported secondary profile pulls back along
the term-context transport to a valuated secondary reading at the actual index.
The resulting actual reading still carries the fiber-local topic cover.
\end{theorem}

\begin{theorem}[Actual pullback readings recover valuations]\label{thm:claim-secondary-pullback-valuation-actual}\lean{Obsmetrologia.Claim.IntensionalFrame.secondaryActualPullbackValuation_at_actual}\leanok\uses{thm:claim-secondary-pullback-valuated-reading-cover, thm:claim-anchor-actual-pullback-valuation}
When the pullback target is the anchor's actual index, the pullback-valuated
secondary reading carries the original valuation.
\end{theorem}

\begin{proposition}[Pulled-back secondary readings induce relative patch families]\label{thm:claim-secondary-pullback-patch-family}\lean{Obsmetrologia.Claim.IntensionalFrame.secondaryActualPullbackPatchFamily}\leanok\uses{thm:claim-secondary-pullback-valuated-reading-cover, thm:claim-valuated-reading-topic-patch-family}
After target valuation support is pulled back along the term-context transport,
the actual secondary reading also induces a relative patch family over its
valuation-read support.
\end{proposition}

\begin{theorem}[Actual pullback patch supports recover secondary supports]\label{thm:claim-secondary-pullback-patch-support-actual}\lean{Obsmetrologia.Claim.IntensionalFrame.secondaryActualPullbackPatchFamily_support_at_actual}\leanok\uses{thm:claim-secondary-pullback-patch-family, thm:claim-secondary-patch-family, thm:claim-secondary-pullback-valuation-actual}
At the actual index, the pulled-back patch family's support mask is the ordinary
secondary patch-family support mask.
\end{theorem}

\begin{theorem}[Actual pullback patch lists recover secondary patch lists]\label{thm:claim-secondary-pullback-patch-list-actual}\lean{Obsmetrologia.Claim.IntensionalFrame.secondaryActualPullbackPatchFamily_patches_at_actual}\leanok\uses{thm:claim-secondary-pullback-patch-support-actual}
At the actual index, the pulled-back patch family has the same patch list as the
ordinary secondary patch family.
\end{theorem}

\begin{theorem}[Pulled-back secondary patch families cover their support]\label{thm:claim-secondary-pullback-patch-family-covers-support}\lean{Obsmetrologia.Claim.IntensionalFrame.secondaryActualPullbackPatchFamily_covers_support}\leanok\uses{thm:claim-secondary-pullback-patch-family}
The pulled-back secondary patch family covers the support mask read through the
pulled-back valuation.
\end{theorem}

\begin{proposition}[Pulled-back secondary patches promote under full support]\label{thm:claim-secondary-pullback-sheaf-cover}\lean{Obsmetrologia.Claim.IntensionalFrame.secondaryActualPullbackSheafCover}\leanok\uses{thm:claim-secondary-pullback-patch-family-covers-support, thm:claim-valuated-reading-topic-sheaf-cover}
The actual secondary patch family becomes a settling sheaf cover only with a
full-support hypothesis for the pulled-back valuation.
\end{proposition}

\begin{theorem}[Actual pullback full support is ordinary full support]\label{thm:claim-secondary-pullback-full-support-actual}\lean{Obsmetrologia.Claim.IntensionalFrame.actualFullSupport_of_pullback}\leanok\uses{thm:claim-anchor-actual-pullback-valuation}
At the actual index, full support for the pulled-back valuation is full support
for the original valuation.
\end{theorem}

\begin{theorem}[Actual pullback sheaf covers recover secondary patch lists]\label{thm:claim-secondary-pullback-sheaf-patches-actual}\lean{Obsmetrologia.Claim.IntensionalFrame.secondaryActualPullbackSheafCover_patches_at_actual}\leanok\uses{thm:claim-secondary-pullback-sheaf-cover, thm:claim-secondary-sheaf-cover, thm:claim-secondary-pullback-full-support-actual}
At the actual index, the pulled-back secondary sheaf cover has the same patch
list as the ordinary secondary sheaf cover whose full-support hypothesis has
been transported back from the pulled-back valuation.
\end{theorem}

\begin{theorem}[Covers are fiber-local, not temporal]\label{thm:no-temporal-covers}\lean{Obsmetrologia.Site.no_temporal_covers}\leanok\uses{thm:bitopological-site-grothendieck}
A cover refines the clopen fiber at fixed predicate and stratum.  Pure temporal
advancement is carried by morphisms, not by the topology.  Cross-stratum
content must therefore enter through morphisms and later direct limits.
\end{theorem}

\begin{theorem}[The strict site has finite clopen fiber]\label{thm:strict-site-finite-fiber}\lean{Obsmetrologia.Site.strict_site_finite_fiber}\leanok\uses{thm:no-temporal-covers, thm:spine-is-strict-regime}
In the strict regime, the clopen fiber is the four-element Boolean algebra.
This is the finite site on which the first sheaf-cohomology calculation lives.
\end{theorem}
